
\textbf{Definitions and assumptions.}
The QUBE-Servo 2 rotary pendulum comprises a motor-driven horizontal arm
and a pendulum hinged at its tip. Encoders measure the two angles
\cite{qube_manual}. Let $\theta$ denote arm angle and $\alpha$ pendulum
angle from the upward vertical, with $\alpha=0$ upright. Use radians.
Let $\boldsymbol e_r$ point along the arm, $\boldsymbol e_t$ point in the
direction of increasing $\theta$ and $\boldsymbol e_z$ point upward.
Positive $\alpha$ tilts the pendulum towards $-\boldsymbol e_t$.
This explicitly fixes the coupling signs.

The arm and pendulum are approximated as rigid uniform slender rods.
The pendulum centre of mass is at $l_p=L_p/2$. Viscous damping represents
friction. Motor inductance, structural flexibility, nonlinear friction
and pendulum longitudinal-axis inertia are neglected. The motor is direct
drive. These are nominal modelling assumptions, not measurements of
the particular laboratory unit.
\begin{equation}
x=\begin{bmatrix}\theta&\alpha&\dot\theta&\dot\alpha\end{bmatrix}^{T},
\qquad u=V_m,\qquad y=\begin{bmatrix}\theta&\alpha\end{bmatrix}^{T}.
\end{equation}
Thus there are $n=4$ states, $m=1$ input and $p=2$ outputs. Velocities
are states even though only the angles are directly measured.

\textbf{Motor electrical dynamics.}
Kirchhoff's voltage law and the motor relations give
\begin{align}
V_m&=R_mi_m+L_m\frac{di_m}{dt}+e_b,\\
e_b&=k_m\dot\theta,\qquad \tau_m=k_ti_m.
\end{align}
Neglecting inductance,
\begin{equation}\label{eq:motor}
i_m=\frac{V_m-k_m\dot\theta}{R_m},\qquad
\boxed{\tau_m=\frac{k_t}{R_m}V_m-\frac{k_tk_m}{R_m}\dot\theta.}
\end{equation}
Voltage generates torque and back emf reduces the current as speed rises.

\textbf{Inertias and energy.}
For a uniform rod, $dm=(m/L)\,dr$. The relevant inertias are
\begin{align}
J_{\mathrm{arm}}&=\int_0^{L_r}r^2\frac{m_r}{L_r}\,dr
=\frac13m_rL_r^2,\qquad J_r=J_{\mathrm{arm}}+J_m+J_h,\\
J_p&=\int_{-L_p/2}^{L_p/2}r^2\frac{m_p}{L_p}\,dr
=\frac1{12}m_pL_p^2.
\end{align}
Here $J_r$ includes the arm, rotor and hub. The arm rod approximation
excludes the separately added rotor and hub. No inertia-disc module is
included. $J_p$ is about the pendulum centre of mass; the hinge inertia is
$J_p+m_pl_p^2=\frac13m_pL_p^2$ by the parallel-axis theorem.

The pendulum centre-of-mass position is
\begin{equation}
\boldsymbol r_p=L_r\boldsymbol e_r-l_p\sin\alpha\,\boldsymbol e_t
+l_p\cos\alpha\,\boldsymbol e_z.
\end{equation}
Using $\dot{\boldsymbol e}_r=\dot\theta\boldsymbol e_t$ and
$\dot{\boldsymbol e}_t=-\dot\theta\boldsymbol e_r$ gives
\begin{align}
\dot{\boldsymbol r}_p&=l_p\sin\alpha\,\dot\theta\,\boldsymbol e_r
+(L_r\dot\theta-l_p\cos\alpha\,\dot\alpha)\boldsymbol e_t
-l_p\sin\alpha\,\dot\alpha\,\boldsymbol e_z,\\
\|\dot{\boldsymbol r}_p\|^2&=(L_r^2+l_p^2\sin^2\alpha)\dot\theta^2
+l_p^2\dot\alpha^2-2L_rl_p\cos\alpha\,\dot\theta\dot\alpha.
\end{align}
Neglecting longitudinal-axis inertia, pendulum rotation about its centre
of mass contributes
\[\frac12J_p(\dot\alpha^2+\sin^2\alpha\,\dot\theta^2).\]
Therefore the total kinetic and potential energies are
\begin{align}
T&=\frac12J_r\dot\theta^2+\frac12m_p\|\dot{\boldsymbol r}_p\|^2
+\frac12J_p(\dot\alpha^2+\sin^2\alpha\,\dot\theta^2),\\
U&=m_pgl_p\cos\alpha.
\end{align}
Potential energy is referenced to the hinge height and is maximal upright.

Define
\begin{equation}\label{eq:constants}
a=J_r+m_pL_r^2,\quad b=J_p+m_pl_p^2,\quad
c=m_pL_rl_p,\quad h=m_pgl_p.
\end{equation}
The Lagrangian is
\begin{equation}\label{eq:lagrangian}
\mathcal L=T-U=
\frac12(a+b\sin^2\alpha)\dot\theta^2+\frac12b\dot\alpha^2
-c\cos\alpha\,\dot\theta\dot\alpha-h\cos\alpha.
\end{equation}

\textbf{Euler--Lagrange equations.}
For $q\in\{\theta,\alpha\}$,
\begin{equation}
\frac{d}{dt}\left(\frac{\partial\mathcal L}{\partial\dot q}\right)
-\frac{\partial\mathcal L}{\partial q}=Q_q,\qquad
Q_\theta=\tau_m-D_r\dot\theta,\quad Q_\alpha=-D_p\dot\alpha.
\end{equation}
The arm derivatives are
\begin{align}
\frac{\partial\mathcal L}{\partial\theta}&=0,\\
\frac{\partial\mathcal L}{\partial\dot\theta}
&=(a+b\sin^2\alpha)\dot\theta-c\cos\alpha\,\dot\alpha,\\
\frac{d}{dt}\left(\frac{\partial\mathcal L}{\partial\dot\theta}\right)
&=(a+b\sin^2\alpha)\ddot\theta
+2b\sin\alpha\cos\alpha\,\dot\alpha\dot\theta
-c\cos\alpha\,\ddot\alpha+c\sin\alpha\,\dot\alpha^2.
\end{align}
The pendulum derivatives are
\begin{align}
\frac{\partial\mathcal L}{\partial\dot\alpha}
&=b\dot\alpha-c\cos\alpha\,\dot\theta,\\
\frac{d}{dt}\left(\frac{\partial\mathcal L}{\partial\dot\alpha}\right)
&=b\ddot\alpha-c\cos\alpha\,\ddot\theta
+c\sin\alpha\,\dot\alpha\dot\theta,\\
\frac{\partial\mathcal L}{\partial\alpha}
&=b\sin\alpha\cos\alpha\,\dot\theta^2
+c\sin\alpha\,\dot\theta\dot\alpha+h\sin\alpha.
\end{align}
Define $k_v=k_t/R_m$ and $d=D_r+k_tk_m/R_m$. The resulting
nonlinear mechanical equations, including motor torque, are
\begin{samepage}
\begin{subequations}\label{eq:eom}
\begin{align}
(a+b\sin^2\alpha)\ddot\theta-c\cos\alpha\,\ddot\alpha
+2b\sin\alpha\cos\alpha\,\dot\alpha\dot\theta
+c\sin\alpha\,\dot\alpha^2&=k_vu-d\dot\theta,\\
-c\cos\alpha\,\ddot\theta+b\ddot\alpha
-b\sin\alpha\cos\alpha\,\dot\theta^2-h\sin\alpha&=-D_p\dot\alpha.
\end{align}
\end{subequations}
\end{samepage}
Fixing the arm and neglecting friction gives $b\ddot\alpha=h\sin\alpha$.
A small positive tilt therefore accelerates away from upright, which
provides a physical check on the gravity sign.

\textbf{Nonlinear state-space representation.}
Define the torque expressions
\begin{align}
r_1(x,u)&=k_vu-dx_3-2b\sin x_2\cos x_2\,x_4x_3-c\sin x_2\,x_4^2,\\
r_2(x,u)&=b\sin x_2\cos x_2\,x_3^2+h\sin x_2-D_px_4.
\end{align}
The equations of motion become
\begin{equation}
\underbrace{\begin{bmatrix}
a+b\sin^2\alpha&-c\cos\alpha\\-c\cos\alpha&b
\end{bmatrix}}_{M(\alpha)}
\begin{bmatrix}\ddot\theta\\\ddot\alpha\end{bmatrix}
=\begin{bmatrix}r_1\\r_2\end{bmatrix}.
\end{equation}
The mass-matrix determinant and inverse are
\begin{align}
\Delta(\alpha)&=b(a+b\sin^2\alpha)-c^2\cos^2\alpha,\\
M^{-1}(\alpha)&=\frac1{\Delta(\alpha)}
\begin{bmatrix}b&c\cos\alpha\\c\cos\alpha&a+b\sin^2\alpha\end{bmatrix}.
\end{align}
Hence
\begin{equation}\label{eq:nonlinear_ss}
\boxed{\dot x=f(x,u)=
\begin{bmatrix}
x_3\\x_4\\
(br_1+c\cos x_2\,r_2)/\Delta(x_2)\\
(c\cos x_2\,r_1+(a+b\sin^2x_2)r_2)/\Delta(x_2)
\end{bmatrix}.}
\end{equation}
For the parameters in Section~\ref{sec:model}, $ab-c^2>0$ and
$\Delta(\alpha)=(ab-c^2)+(b^2+c^2)\sin^2\alpha>0$, so this idealised
model is nonsingular for every pendulum angle.

\textbf{Equilibrium and Jacobian linearisation.}
An upright stationary equilibrium is
\begin{equation}
x_e=\begin{bmatrix}\theta_e&0&0&0\end{bmatrix}^{T},\qquad u_e=0.
\end{equation}
Any fixed $\theta_e$ is an ideal equilibrium, since no restoring torque
depends on arm angle. Choose $\theta_e=0$ for the numerical model.
Define $\delta x=x-x_e$, $\delta u=u-u_e$ and $\delta y=y-y_e$.
The first-order Taylor expansion is
\begin{equation}\label{eq:ct_ss}
\delta\dot x=A_c\delta x+B_c\delta u,\qquad
\delta y=C_c\delta x+D_c\delta u,
\end{equation}
with Jacobians
\begin{equation}
A_c=\left.\frac{\partial f}{\partial x}\right|_e,\qquad
B_c=\left.\frac{\partial f}{\partial u}\right|_e.
\end{equation}
At equilibrium, $r_1=r_2=0$ and $\Delta_0=ab-c^2$. Also,
\begin{equation}
\left.\frac{\partial(r_1,r_2)}{\partial x}\right|_e
=\begin{bmatrix}0&0&-d&0\\0&h&0&-D_p\end{bmatrix},\qquad
\left.\frac{\partial(r_1,r_2)}{\partial u}\right|_e
=\begin{bmatrix}k_v\\0\end{bmatrix}.
\end{equation}
Derivatives of $M^{-1}$ multiply the zero equilibrium torque vector.
Multiplying the above derivatives by
$M^{-1}(0)=\Delta_0^{-1}\begin{bmatrix}b&c\\c&a\end{bmatrix}$ gives
\begin{equation}\label{eq:Ac_symbolic}
\boxed{A_c=
\begin{bmatrix}
0&0&1&0\\0&0&0&1\\
0&ch/\Delta_0&-bd/\Delta_0&-cD_p/\Delta_0\\
0&ah/\Delta_0&-cd/\Delta_0&-aD_p/\Delta_0
\end{bmatrix},\quad
B_c=\begin{bmatrix}0\\0\\bk_v/\Delta_0\\ck_v/\Delta_0\end{bmatrix}.}
\end{equation}
Since the outputs are the two angles with no direct voltage feedthrough,
\begin{equation}\label{eq:CD}
\boxed{C_c=\begin{bmatrix}1&0&0&0\\0&1&0&0\end{bmatrix},
\qquad D_c=\begin{bmatrix}0\\0\end{bmatrix}.}
\end{equation}
This linearisation describes small deviations about upright; large-angle
motion requires the nonlinear model in~\eqref{eq:nonlinear_ss}.

\subsection{System parameters}

Table~\ref{tab:params} lists the manufacturer-supplied parameters relevant
to the pendulum configuration. The inertia-disc parameters are excluded.
Motor inductance is listed to justify its neglect, rather than being
retained as a model state.
\begin{table}[H]\centering
\caption{Parameters from Table 2.2 of the QUBE-Servo 2 user manual~\cite{qube_manual}.}
\label{tab:params}
\begin{tabular}{lll}
\toprule
\textbf{Symbol}&\textbf{Description}&\textbf{Value}\\
\midrule
$R_m$&Motor resistance&$8.4~\Omega$\\
$k_t$&Torque constant&$0.042~\mathrm{N\,m/A}$\\
$k_m$&Back-emf constant&$0.042~\mathrm{V\,s/rad}$\\
$J_m$&Motor rotor inertia&$4.0\times10^{-6}~\mathrm{kg\,m^2}$\\
$J_h$&Attachment hub inertia&$0.6\times10^{-6}~\mathrm{kg\,m^2}$\\
$L_m$&Inductance (neglected)&$1.16\times10^{-3}~\mathrm H$\\
$m_r$&Arm mass&$0.095~\mathrm{kg}$\\
$L_r$&Arm length, motor pivot to tip&$0.085~\mathrm m$\\
$m_p$&Pendulum mass&$0.024~\mathrm{kg}$\\
$L_p$&Pendulum full length&$0.129~\mathrm m$\\
&Encoder resolution (quadrature)&$2\pi/2048~\mathrm{rad/count}$\\
&Recommended amplifier output&$\pm10~\mathrm V$\\
\bottomrule
\end{tabular}
\end{table}
The motor electrical time constant is
\begin{equation}
\tau_{\mathrm{elec}}=\frac{L_m}{R_m}
=\frac{1.16\times10^{-3}}{8.4}
=1.3810\times10^{-4}~\mathrm s.
\end{equation}
This is short compared with the mechanical time scales, motivating the
quasi-static current approximation. Retaining inductance would introduce
motor current as a fifth state.

The damping values in Table~\ref{tab:additional} are adopted from the
manufacturer's student workbook, \emph{Pendulum State-Space Model}, in-lab
exercise 6~\cite{qube_courseware}. That source reports them as experimentally
obtained example values; they were not identified from this group's
specific laboratory unit. No new experimental identification is claimed.
\begin{table}[H]\centering
\caption{Additional constants and their provenance.}
\label{tab:additional}
\begin{tabular}{llll}
\toprule
\textbf{Symbol}&\textbf{Description}&\textbf{Value}&\textbf{Basis}\\
\midrule
$D_r$&Arm damping&$0.0015~\mathrm{N\,m\,s/rad}$&Workbook example\\
$D_p$&Pendulum damping&$0.0005~\mathrm{N\,m\,s/rad}$&Workbook example\\
$g$&Gravity&$9.81~\mathrm{m/s^2}$&Standard approximation\\
\bottomrule
\end{tabular}
\end{table}

\begin{samepage}
The derived dimensions and inertias are
\begin{align}
l_p&=L_p/2=0.0645~\mathrm m,\\
J_{\mathrm{arm}}&=\tfrac13(0.095)(0.085)^2
=2.2879167\times10^{-4}~\mathrm{kg\,m^2},\\
J_r&=J_{\mathrm{arm}}+4.0\times10^{-6}+0.6\times10^{-6}
=2.3339167\times10^{-4}~\mathrm{kg\,m^2},\\
J_p&=\tfrac1{12}(0.024)(0.129)^2
=3.3282\times10^{-5}~\mathrm{kg\,m^2}.
\end{align}
\end{samepage}
The constants used in the matrices are
\begin{align}
a&=4.0679167\times10^{-4}~\mathrm{kg\,m^2},&
b&=1.33128\times10^{-4}~\mathrm{kg\,m^2},\\
c&=1.3158\times10^{-4}~\mathrm{kg\,m^2},&
h&=0.01518588~\mathrm{N\,m},\\
d&=0.00171~\mathrm{N\,m\,s/rad},&
k_v&=0.005~\mathrm{N\,m/V},\\
\Delta_0&=3.68420646\times10^{-8}~\mathrm{kg^2\,m^4}.&&
\end{align}
Substituting into~\eqref{eq:Ac_symbolic} gives
\begin{equation}\label{eq:Ac_numeric}
A_c\simeq
\begin{bmatrix}
0&0&1&0\\0&0&0&1\\
0&54.2357795&-6.17904785&-1.78573054\\
0&167.674898&-6.10719845&-5.52075014
\end{bmatrix},\quad
B_c\simeq\begin{bmatrix}0\\0\\18.0673914\\17.8573054\end{bmatrix}.
\end{equation}
The matrices $C_c,D_c$ remain exactly those in~\eqref{eq:CD}.
These are nominal simulation parameters. Actual mass distribution, damping,
encoder calibration and unmodelled friction may produce differences from
hardware measurements.

\subsection{Discretisation and sampling time}

Choose the sampling time
\begin{equation}
\boxed{T_s=0.001~\mathrm s,\qquad f_s=1/T_s=1000~\mathrm{Hz}.}
\end{equation}
This matches the value in the supplied ControLab template's
\texttt{Hardware\_settings.m}~\cite{hardware_template}. It is also short
relative to the mechanical dynamics, since
$T_s\max_i|\lambda_i(A_c)|\simeq0.01763\ll1$ for this model. Hardware execution timing
and achieved closed-loop performance must still be checked during
implementation.

Zero-order hold (ZOH) assumes that motor voltage is held constant between
controller updates. The exact discretisation of the linearised model
under this assumption is
\begin{equation}\label{eq:zoh}
A=e^{A_cT_s},\qquad
B=\int_0^{T_s}e^{A_c\tau}B_c\,d\tau,\qquad C=C_c,\quad D=D_c.
\end{equation}
The MATLAB function \texttt{c2d} evaluates these expressions without
requiring an inverse of $A_c$. This matters because the arm-position mode
makes $A_c$ singular.
The discrete-time model is
\begin{subequations}\label{eq:dt_ss}
\begin{align}
\delta x[k+1]&=A\delta x[k]+B\delta u[k],\\
\delta y[k]&=C\delta x[k]+D\delta u[k].
\end{align}
\end{subequations}
At the chosen zero equilibrium the deviation variables equal the
original variables. The calculated matrices are
\begin{equation}\label{eq:A_discrete}
A\simeq
\begin{bmatrix}
1&2.70127657\!\times\!10^{-5}&9.96918625\!\times\!10^{-4}&-8.80391044\!\times\!10^{-7}\\
0&1.00008362958&-3.04176177\!\times\!10^{-6}&9.97274363\!\times\!10^{-4}\\
0&0.05392103923&0.99384536883&-0.00174835485\\
0&0.16705290505&-0.00607175726&0.99458335877
\end{bmatrix},
\end{equation}
\begin{equation}
B\simeq\begin{bmatrix}
9.00986699\times10^{-6}\\8.89404027\times10^{-6}\\
0.0179959975588\\0.0177536761881
\end{bmatrix},\quad
C=\begin{bmatrix}1&0&0&0\\0&1&0&0\end{bmatrix},\quad
D=\begin{bmatrix}0\\0\end{bmatrix}.
\end{equation}
The dimensions are $A\in\mathbb R^{4\times4}$,
$B\in\mathbb R^{4\times1}$, $C\in\mathbb R^{2\times4}$ and
$D\in\mathbb R^{2\times1}$. Displayed results are rounded;
full precision must be retained in MATLAB for controller and observer
design.

